\documentclass[12pt]{article}
\usepackage[T1]{fontenc}
\usepackage[utf8]{inputenc}
\usepackage{lmodern}
\usepackage{textcomp}
\usepackage{enumitem}
\usepackage{geometry}
\geometry{margin=1in}
\begin{document}
\begin{center}
Answer Key - Religious Studies - K-12 Level Assessment 25 \
\small (Answers are ordered exactly as in the question paper.)
\end{center}

\section*{Multiple Choice}
\begin{enumerate}[label=\arabic*.]
\item \textbf{Answer:} D
\\ \textit{Options:} A) Asceticism; B) Bodily pleasure; C) Music; D) Sexual union
\\ \textit{Note:} Sexual union is considered a central component of tantric Buddhism as it symbolizes the integration of male and female energies, which is essential for achieving spiritual enlightenment and union with the divine.
\item \textbf{Answer:} C
\\ \textit{Options:} A) Toward the end of the first millennium BCE; B) Towards the middle of the first millennium BCE; C) Toward the end of the second millennium BCE; D) Toward the middle of the second millennium BCE
\\ \textit{Note:} The correct answer is toward the end of the second millennium BCE because this period marks the time when the various versions of the Epic of Gilgamesh were standardized and compiled into a cohesive narrative, reflecting the cultural and literary developments of ancient Mesopotamia.
\item \textbf{Answer:} D
\\ \textit{Options:} A) Amulets; B) Meditation; C) Rituals; D) Grand festivals
\\ \textit{Note:} The term "matsuri" in Japanese culture specifically refers to traditional grand festivals that celebrate and honor deities, seasonal changes, or significant events, thus making "grand festivals" the accurate definition.
\item \textbf{Answer:} A
\\ \textit{Options:} A) Julian of Norwich; B) Catherine of Siena; C) Teresa of Avila; D) John of the Cross
\\ \textit{Note:} Julian of Norwich, a 14 th-century mystic, articulated the belief that evil serves as a distortion that ultimately illuminates and enhances the understanding of divine love, emphasizing God's mercy and the potential for redemption.
\item \textbf{Answer:} B
\\ \textit{Options:} A) Theravada; B) Chan-Zen; C) Pure Land; D) Yogicara
\\ \textit{Note:} The "Flower Sermon," in which the Buddha silently held up a flower to convey enlightenment, is a foundational teaching of Chan-Zen Buddhism, emphasizing direct experience and intuition over intellectual understanding.
\end{enumerate}

\section*{True / False}
\begin{enumerate}[label=\arabic*.]
\setcounter{enumi}{5}
\item \textbf{Answer:} False
\\ \textit{Options:} A) True; B) False
\\ \textit{Note:} The statement is false because Islamic jurisprudence is derived from multiple sources, including the Quran, the Sunnah (traditions of the Prophet Muhammad), Ijtihad (independent reasoning), and consensus (Ijma), not just consensus and personal reasoning.
\item \textbf{Answer:} True
\\ \textit{Options:} A) True; B) False
\\ \textit{Note:} Wearing the hijab serves as a significant expression of Islamic identity for many individuals, symbolizing their faith, cultural heritage, and commitment to religious values in contemporary society.
\item \textbf{Answer:} True
\\ \textit{Options:} A) True; B) False
\\ \textit{Note:} Pope Gregory promoted the cults of St. Peter and St. Paul to strengthen the papal authority by aligning his leadership with the revered apostles, who were foundational figures in the Christian faith and symbolized the legitimacy of the Roman Church.
\item \textbf{Answer:} True
\\ \textit{Options:} A) True; B) False
\\ \textit{Note:} The Bodu Bala Sena is considered an activist and militant organization because it actively promotes Buddhist nationalism and has been involved in aggressive campaigns against religious minorities in Sri Lanka.
\item \textbf{Answer:} True
\\ \textit{Options:} A) True; B) False
\\ \textit{Note:} The Great Cloud Sutra indeed prophesies the future arrival of Maitreya, who is revered in Buddhist tradition as the next Buddha destined to bring enlightenment to the world.
\end{enumerate}

\section*{Long Form Response}
\begin{enumerate}[label=\arabic*.]
\setcounter{enumi}{10}
\item \textbf{Answer:} Tafsir is the Arabic term for commentary on the Qur'an, and it plays a crucial role in Islamic studies by helping readers understand the meaning and context of the sacred text. Through tafsir, scholars interpret the verses of the Qur'an, explaining their linguistic, historical, and theological significance. This process not only clarifies the message of the Qur'an but also addresses various interpretations that arise due to cultural and temporal differences. By providing insights into the text, tafsir enhances the understanding of Islamic teachings and values, making the Qur'an more accessible to believers and scholars alike. Ultimately, tafsir serves as a bridge between the divine message and human comprehension, enriching the spiritual and intellectual lives of Muslims.
\\ \textit{Note:} correct because it accurately defines tafsir as an essential tool for interpreting the Qur'an, emphasizing its role in elucidating the text's meaning, context, and significance while addressing diverse interpretations shaped by cultural and temporal factors.
\item \textbf{Answer:} In the early church, deaconesses played several important roles that significantly contributed to the community and religious practices. They were often responsible for assisting with the baptism of women, providing care for the sick, and supporting those in need, which highlighted their role in fostering compassion and community support. Additionally, deaconesses acted as a bridge between the clergy and the congregation, ensuring that the needs of the women and families were communicated effectively. Their contributions were vital in establishing a more inclusive church environment, emphasizing the importance of service and ministry in early Christian life. Overall, the various roles of deaconesses underscored their significance in promoting both spiritual and social welfare within the community.
\\ \textit{Note:} correct because it identifies the multifaceted roles of deaconesses in the early church, emphasizing their contributions to women's baptism, community care, and their vital communicative role between clergy and congregation, which collectively illustrate their significance in religious practices and community support.
\end{enumerate}

\vfill
\noindent\textit{End of Answer Key}
\end{document}